\documentclass[UTF8]{ctexart}
\usepackage[a4paper,landscape,margin=4mm]{geometry}
\usepackage{multicol}
\usepackage{enumitem}
\usepackage{xcolor}
\usepackage{amsmath,amssymb}
\usepackage{listings}
\usepackage{titlesec}

\pagestyle{empty}
\setlength{\parindent}{0pt}
\setlength{\parskip}{0pt}
\setlength{\columnsep}{3.2mm}
\setlength{\columnseprule}{0.08pt}
\setlength{\premulticols}{0pt}
\setlength{\postmulticols}{0pt}
\setlength{\multicolsep}{0pt}
\linespread{0.78}
\titleformat{\section}{\bfseries\fontsize{8.7pt}{8.9pt}\selectfont}{}{0pt}{}
\titlespacing*{\section}{0pt}{1pt}{0.5pt}
\setlist[itemize]{leftmargin=1.0em,nosep,topsep=0pt,partopsep=0pt,itemsep=0pt,parsep=0pt}
\lstset{
  basicstyle=\ttfamily\fontsize{5.6pt}{6.0pt}\selectfont,
  columns=fullflexible,
  keepspaces=true,
  breaklines=true,
  showstringspaces=false,
  aboveskip=0.5pt,
  belowskip=0.5pt,
  xleftmargin=0pt,
  frame=none
}
\newcommand{\blk}[1]{\vspace{0.8pt}\noindent\colorbox{black!10}{\parbox{\dimexpr\linewidth-2\fboxsep}{\bfseries #1}}\vspace{0.6pt}}
\newcommand{\tri}{$\Rightarrow$}
\newcommand{\kw}[1]{\textbf{#1}}
\newcommand{\tinyline}{\vspace{0.4pt}\hrule height 0.08pt\vspace{0.4pt}}

\begin{document}
\fontsize{8.55pt}{9.0pt}\selectfont

\begin{multicols*}{4}
\blk{总知识图谱}
现实世界 \tri 信息世界 \tri 机器世界。业务需求经 E-R 抽象为实体、属性、联系，再转关系模式；关系模式经函数依赖和范式消除冗余；SQL 在 DBMS 上完成定义、操纵、控制；DBMS 通过索引、执行计划、事务、并发控制保证效率和正确性；分布式和 NoSQL 是规模、可用性、模型灵活性的扩展。

数据库系统主线：\kw{数据模型}定结构和约束；\kw{数据库设计}定模式；\kw{SQL}表达查询和更新；\kw{存储结构/索引}定访问路径；\kw{查询执行/优化}定代价；\kw{事务/锁}定并发正确性；\kw{复制/分片/CAP}定分布式取舍。

\blk{1 数据库系统概论}
\kw{数据}是描述客观事物的符号；\kw{信息}是加工解释后的有用数据；\kw{数据处理}包括采集、组织、存储、加工、传播。人工管理：程序和数据强耦合、无共享。文件系统：有文件独立性但冗余、不一致、共享差。数据库系统：统一管理、共享高、冗余低、独立性高、完整性和安全性强、并发与恢复由 DBMS 管理。

DB：长期存储、有组织、可共享的数据集合。DBMS：位于用户和 OS 之间，负责定义、操纵、运行管理、建立维护。DBS=DB+DBMS+应用+DBA+用户。DBA：模式定义、权限安全、性能监控、备份恢复、完整性维护。

\kw{三级模式}：外模式=用户视图；模式=全局逻辑结构；内模式=物理存储。两级映像：外模式/模式保证逻辑独立性；模式/内模式保证物理独立性。数据模型三要素：数据结构、数据操作、完整性约束。常见逻辑模型：层次、网状、关系、对象关系、NoSQL。C/S 客户端重，B/S 浏览器轻；云数据库关注弹性、托管、按需付费。

\blk{2 关系模型和关系代数}
关系=二维表；元组=行；属性=列；域=取值范围；关系模式 $R(A_1,\ldots,A_n)$；关系实例是某一时刻的元组集合。码：超码唯一标识；候选码最小超码；主码被选作主标识；外码引用其他关系主码或候选码。关系性质：列同质、列名唯一、行列无序、任意两个元组不同、属性值原子。

传统集合运算要求并兼容：并、差、交、笛卡尔积。专门关系运算：选择 $\sigma_F(R)$ 取行；投影 $\pi_A(R)$ 取列并去重；连接 $\theta$-join；自然连接同名属性等值且只保留一份；除法用于“满足所有”。

常用等价：$R\cap S=R-(R-S)$；$R\bowtie_F S=\sigma_F(R\times S)$；自然连接=等值连接+投影去重同名列。除法 $R(XY)\div S(Y)$ 结果为满足对每个 $y\in S$ 均有 $(x,y)\in R$ 的 $x$。

\blk{函数依赖和范式}
$X\to Y$：任意两元组若 $X$ 相同则 $Y$ 相同。平凡依赖：$Y\subseteq X$。完全依赖：$X\to Y$ 且任意真子集不决定 $Y$；部分依赖反之。传递依赖：$X\to Y,Y\to Z,Y\nrightarrow X$ 则 $Z$ 传递依赖于 $X$。闭包 $X^+$：由 $X$ 在依赖集 $F$ 下可推出的全部属性；若 $X^+$ 含全部属性且 $X$ 最小，则 $X$ 是候选码。

Armstrong：自反、增广、传递；常用推出：合并、分解、伪传递。最小覆盖：右部单属性，左部无多余属性，依赖无冗余。

范式速判：1NF 属性原子；2NF 在 1NF 上消除非主属性对候选码的部分依赖；3NF 消除非主属性对候选码的传递依赖，等价为对每个非平凡 $X\to A$，$X$ 是超码或 $A$ 是主属性；BCNF 要求每个非平凡决定因素 $X$ 都是超码。强度：BCNF \tri 3NF \tri 2NF \tri 1NF。异常：插入异常、删除异常、修改异常来自冗余和依赖混杂。

分解：无损连接要求分解后自然连接能还原原关系；二元分解 $R\to R_1,R_2$ 无损当且仅当 $(R_1\cap R_2)\to R_1$ 或 $(R_1\cap R_2)\to R_2$。保持依赖：各子模式依赖投影的闭包能推出原 $F$。3NF 常可兼顾无损和保持依赖；BCNF 可能牺牲保持依赖。

\blk{完整性}
实体完整性：主码非空且唯一。参照完整性：外码为空或等于被参照关系某主码值；删除/更新可选拒绝、级联、置空、置默认。用户自定义完整性：类型、范围、唯一、检查、默认、业务规则。

\blk{3 数据库设计}
设计阶段：需求分析 \tri 概念结构 \tri 逻辑结构 \tri 物理结构 \tri 行为设计 \tri 实施 \tri 运行维护。需求分析产物：数据流图、数据字典、需求说明；关注数据需求、处理需求、安全完整性、性能需求。

DFD：外部实体、处理、数据流、数据存储；逐层分解，父子图输入输出平衡。数据字典记录数据项、数据结构、数据流、数据存储、处理过程，是后续设计依据。

概念设计：自顶向下、自底向上、逐步扩张、混合策略。E-R：实体矩形，属性椭圆，联系菱形；联系基数 1:1、1:n、m:n；弱实体依赖强实体；属性可简单、复合、多值、派生。

E-R 转关系：实体 \tri 关系，主键为实体标识；多值属性 \tri 新关系，含实体键+属性值；1:1 可任选一方加入对方键和联系属性，或独立关系；1:n 在 n 方加入 1 方键和联系属性；m:n 建独立联系关系，键为各实体键组合并含联系属性；多元联系一般建独立关系；自联系按基数处理，角色名区分外键。

逻辑设计：把 E-R 转关系后按函数依赖规范化，检查码、外码、完整性和查询需求。物理设计：存储结构、索引、聚簇、分区、反规范化、访问路径；原则是高频选择/连接/排序分组字段优先，低选择性或频繁更新字段慎建索引。行为设计：事务、过程、触发器、报表、权限、备份恢复流程。实施维护：建库建表、装载数据、调试应用、监控、重组、备份、升级。

\blk{4 MySQL 总览}
MySQL 体系：连接层、服务层、存储引擎层、文件系统。执行流程：连接认证 \tri 解析 SQL \tri 预处理/权限检查 \tri 优化器生成计划 \tri 执行器调用引擎 \tri 返回结果。Server 管逻辑和优化；存储引擎管表数据、索引、事务能力。

SQL 分类：DDL 定义库表视图索引；DML 插入修改删除；DQL 查询；DCL 授权撤销；TCL 提交回滚。语句以分号结束，关键字不区分大小写，字符串用引号，标识符冲突可用反引号。

\begin{lstlisting}
mysql -h host -P 3306 -u user -p
SHOW DATABASES; USE db; SHOW TABLES;
SHOW VARIABLES LIKE 'character_set%';
SHOW ENGINES; SHOW PROCESSLIST;
\end{lstlisting}

\blk{5 库表管理}
引擎：InnoDB 支持事务、行锁、外键、崩溃恢复，默认首选；MyISAM 表锁、无事务无外键、读多写少；MEMORY 数据在内存、重启丢失；CSV 文本；ARCHIVE 压缩归档；BLACKHOLE 丢弃写入。B+Tree 适合范围和排序；Hash 等值快但不支持范围、排序、最左前缀。

字符集决定字符如何编码，校对规则决定比较排序。推荐 utf8mb4。层级：服务器 \tri 数据库 \tri 表 \tri 列，低层未指定继承高层。

数据类型：整数 TINYINT/SMALLINT/INT/BIGINT，可 UNSIGNED；小数 DECIMAL 精确，FLOAT/DOUBLE 近似；日期 DATE/TIME/DATETIME/TIMESTAMP/YEAR；字符串 CHAR 定长、VARCHAR 变长、TEXT 大文本、BLOB 二进制、ENUM 单选、SET 多选。金额用 DECIMAL，身份证电话用字符串。

\begin{lstlisting}
CREATE DATABASE db DEFAULT CHARSET utf8mb4;
DROP DATABASE db;
CREATE TABLE t(
  id INT PRIMARY KEY AUTO_INCREMENT,
  name VARCHAR(50) NOT NULL,
  dept_id INT,
  salary DECIMAL(10,2) DEFAULT 0,
  UNIQUE(name),
  CHECK(salary>=0),
  FOREIGN KEY(dept_id) REFERENCES dept(id)
) ENGINE=InnoDB DEFAULT CHARSET=utf8mb4;
DESC t; SHOW CREATE TABLE t;
ALTER TABLE t ADD col INT;
ALTER TABLE t MODIFY col BIGINT NOT NULL;
ALTER TABLE t CHANGE col new_col INT;
ALTER TABLE t DROP COLUMN col;
RENAME TABLE old TO new; DROP TABLE t;
\end{lstlisting}

约束：PRIMARY KEY 唯一且非空；UNIQUE 唯一可允许空；NOT NULL；DEFAULT；CHECK；FOREIGN KEY。外键要求类型长度兼容、被参照列有索引，InnoDB 支持。

\blk{6 数据操纵和查询}
\begin{lstlisting}
INSERT INTO t(a,b) VALUES(1,'x'),(2,'y');
INSERT INTO t SET a=1,b='x';
INSERT INTO t2(a,b) SELECT a,b FROM t;
REPLACE INTO t(id,a) VALUES(1,'new');
UPDATE t SET a=a+1 WHERE id=1;
DELETE FROM t WHERE id=1;
TRUNCATE TABLE t;
\end{lstlisting}
DELETE 逐行删、可回滚、触发器可能执行；TRUNCATE 重建表、快、通常重置自增、不带 WHERE；DROP 删除表结构。

\begin{lstlisting}
SELECT [DISTINCT] expr AS alias
FROM t
[JOIN s ON t.k=s.k]
WHERE row_condition
GROUP BY cols
HAVING group_condition
ORDER BY cols ASC|DESC
LIMIT offset,count;
\end{lstlisting}
逻辑顺序：FROM/JOIN \tri WHERE \tri GROUP BY \tri HAVING \tri SELECT \tri DISTINCT \tri ORDER BY \tri LIMIT。WHERE 过滤行，不能直接用聚合；HAVING 过滤组，可用聚合。COUNT(*) 计行；COUNT(col) 忽略 NULL；SUM/AVG 忽略 NULL。LIKE 中 \% 任意串，\_ 单字符；BETWEEN 闭区间；NULL 用 IS NULL。

连接：内连接只保留匹配；左外连接保留左表全部，右表无匹配补 NULL；自连接同一表取不同别名。子查询：IN 适合集合；EXISTS 看是否存在，相关子查询逐行判断；ANY/SOME 与任一比较，ALL 与全部比较。UNION 默认去重，UNION ALL 不去重且更快。SQL 注入防护：参数化，最小权限，避免拼接用户输入。

\begin{lstlisting}
SELECT d.name,COUNT(*) n
FROM emp e JOIN dept d ON e.dept_id=d.id
WHERE e.status='A'
GROUP BY d.id,d.name
HAVING COUNT(*)>3
ORDER BY n DESC;

SELECT * FROM emp e
WHERE EXISTS(SELECT 1 FROM dept d
             WHERE d.id=e.dept_id);
\end{lstlisting}

\blk{7 索引和视图}
索引本质是冗余的有序访问结构，提升查询、连接、排序、分组，代价是空间、维护开销、写入变慢。聚簇索引叶子含整行，表数据按主键组织；非聚簇/二级索引叶子含键和主键，可能回表。覆盖索引：查询列全在索引中，避免回表。

B+Tree 规则：等值前缀+范围；复合索引遵守最左匹配；不能跳过中间列；范围条件之后列通常不能继续用于定位；对列做函数、表达式、隐式类型转换会破坏索引；前导通配 LIKE '\%x' 不能用普通 B+Tree 定位。高选择性、高频 WHERE/JOIN/ORDER/GROUP 列适合建；小表、低选择性、频繁更新列慎建。

\begin{lstlisting}
CREATE INDEX idx_a ON t(a);
CREATE UNIQUE INDEX idx_ab ON t(a,b);
ALTER TABLE t ADD INDEX idx_c(c);
SHOW INDEX FROM t;
DROP INDEX idx_a ON t;

CREATE VIEW v AS SELECT a,b FROM t WHERE c>0;
ALTER VIEW v AS SELECT a FROM t;
DROP VIEW v;
\end{lstlisting}
视图是虚表，存定义不一定存数据；作用：简化查询、封装复杂逻辑、安全隔离、逻辑独立。可更新视图通常要求来自单表、保留键、无聚合/分组/DISTINCT/UNION/复杂表达式。
\blk{8 存储过程、函数、游标}
存储过程：封装一组 SQL，可有 IN/OUT/INOUT 参数，可返回多个结果集，适合批处理和业务逻辑。函数：必须 RETURN 单值，可嵌入表达式，限制更多。共同点：预编译、复用、减少网络往返；缺点是迁移性差、调试困难、过度使用会把业务逻辑锁进数据库。

\begin{lstlisting}
DELIMITER //
CREATE PROCEDURE p(IN x INT, OUT y INT)
BEGIN
  DECLARE v INT DEFAULT 0;
  SELECT COUNT(*) INTO v FROM t WHERE a>x;
  SET y=v;
END//
DELIMITER ;
CALL p(10,@n); SELECT @n;

CREATE FUNCTION f(x INT) RETURNS INT
DETERMINISTIC
BEGIN
  RETURN x*x;
END;
\end{lstlisting}

变量：局部变量 DECLARE，用户变量 @v，会话变量，系统变量 @@var。赋值 SET 或 SELECT ... INTO。条件处理：DECLARE CONTINUE/EXIT HANDLER FOR SQLEXCEPTION/NOT FOUND/SQLSTATE。游标步骤：DECLARE cursor \tri DECLARE handler \tri OPEN \tri FETCH 循环 \tri CLOSE；只读、逐行、常配合 NOT FOUND 标志。

流程：IF 条件 THEN ... ELSEIF ... ELSE END IF；CASE；LOOP + LEAVE/ITERATE；WHILE 先判定；REPEAT 后判定。查看：SHOW PROCEDURE STATUS；SHOW CREATE PROCEDURE；删除 DROP PROCEDURE/FUNCTION。

\blk{常用函数}
字符串：CONCAT、LENGTH/CHAR\_LENGTH、SUBSTRING、LEFT/RIGHT、TRIM、REPLACE、UPPER/LOWER。数学：ABS、ROUND、CEIL、FLOOR、MOD、RAND。日期：NOW、CURDATE、CURTIME、DATE\_ADD、DATEDIFF、YEAR/MONTH/DAY。条件：IF、IFNULL、NULLIF、CASE WHEN。系统：DATABASE、USER、VERSION。加密/摘要：MD5、SHA2。

\blk{9 触发器、事件、权限}
触发器在 INSERT/UPDATE/DELETE 的 BEFORE/AFTER 自动执行；行级触发器对每行触发。OLD 表示旧行，NEW 表示新行：INSERT 只有 NEW，DELETE 只有 OLD，UPDATE 两者都有。适合审计、派生字段、复杂完整性；不适合复杂业务流程，避免递归和隐式副作用。

\begin{lstlisting}
CREATE TRIGGER trg BEFORE UPDATE ON emp
FOR EACH ROW
BEGIN
  SET NEW.updated_at=NOW();
END;
SHOW TRIGGERS; DROP TRIGGER trg;

SET GLOBAL event_scheduler=ON;
CREATE EVENT ev ON SCHEDULE EVERY 1 DAY
DO DELETE FROM log WHERE created_at<NOW()-INTERVAL 30 DAY;
\end{lstlisting}

权限表层级：user 全局，db 数据库，tables\_priv 表，columns\_priv 列，procs\_priv 过程函数。授权原则：最小权限、角色化、及时回收、避免 root 应用连接。

\begin{lstlisting}
CREATE USER 'u'@'host' IDENTIFIED BY 'pwd';
GRANT SELECT,INSERT ON db.t TO 'u'@'host';
REVOKE INSERT ON db.t FROM 'u'@'host';
SHOW GRANTS FOR 'u'@'host';
DROP USER 'u'@'host';
\end{lstlisting}

\blk{事务和并发控制}
事务是一组逻辑操作单元。ACID：原子性全做或全不做；一致性满足约束和业务规则；隔离性并发互不干扰；持久性提交后永久保存。MySQL 默认 autocommit=1；显式事务用 START TRANSACTION、COMMIT、ROLLBACK、SAVEPOINT。

\begin{lstlisting}
START TRANSACTION;
UPDATE account SET bal=bal-100 WHERE id=1;
UPDATE account SET bal=bal+100 WHERE id=2;
COMMIT; -- or ROLLBACK
\end{lstlisting}

并发异常：脏读=读到未提交；不可重复读=同一行两次读值变；幻读=同一条件两次读行数变；丢失修改=覆盖他人更新。隔离级别：READ UNCOMMITTED 可能全有；READ COMMITTED 防脏读；REPEATABLE READ 防不可重复读，InnoDB 配合 MVCC/Next-Key 减少幻读；SERIALIZABLE 最强但并发最低。

调度：串行调度正确但慢；可串行化调度等价于某串行结果；冲突可串行化可用优先图无环判定。锁粒度：表锁开销小冲突大；行锁并发高开销大；意向锁协调表锁和行锁。共享锁 S 可并读，排他锁 X 独占写。两段锁协议：加锁阶段后不可再加锁，解锁阶段后不可再申请新锁，可保证冲突可串行化。死锁：互相等待；处理靠检测回滚、超时、固定访问顺序、短事务、合适索引。

InnoDB 锁：记录锁、间隙锁、Next-Key 锁；索引条件命中才可能精确行锁，无索引可能扩大锁范围。MyISAM 表级锁，读写互斥明显。

\blk{10 查询执行}
B+树原地更新，读和范围查询强，随机写可能引起页分裂。LSM 树写入先到内存 MemTable 和 WAL，刷盘为 SSTable，后台 Compaction 合并；写吞吐高，读需查内存和多层文件，可用 Bloom Filter 降低无效查找；删除用 tombstone。

查询执行流程：SQL \tri 解析树 \tri 逻辑计划 \tri 优化 \tri 物理计划 \tri 执行算子 \tri 结果。执行方式：物化把中间结果写出，简单但 IO 多；流水线边产生边消费，节省中间存储但受阻塞算子限制。阻塞算子：排序、聚合、去重、某些集合操作。

外部排序：内存缓冲 $M$ 页，关系 $B$ 页。Pass0 生成 $\lceil B/M\rceil$ 个有序段；每趟最多归并 $M-1$ 段。两趟条件 $B\le M(M-1)$。趟数 $1+\lceil\log_{M-1}\lceil B/M\rceil\rceil$，完整读写代价约 $2B\times$ 趟数，若最后直接输出可少一次写。

选择执行：无索引全表扫描；主键/唯一索引等值定位；聚簇索引范围扫描；二级索引可能回表；多条件可索引合并但不一定优。投影若不去重可边读边输出；去重要排序或哈希。聚集可排序分组或哈希分组；集合并交差可排序归并或哈希。

连接算法：一趟连接要求较小关系能放内存，代价约 $B_R+B_S$。元组嵌套循环代价 $B_R+T_RB_S$；块嵌套循环 $B_R+\lceil B_R/(M-2)\rceil B_S$，小表作外表。排序归并：Sort(R)+Sort(S)+$B_R+B_S$，适合已排序或范围输出。哈希连接：分区+探测，通常约 $3(B_R+B_S)$，适合等值连接。索引连接：外表每元组用内表索引查找，适合外表小且内表连接列有索引。

\blk{11 查询优化}
优化目标是在等价结果下选低代价计划。逻辑优化改写关系代数；物理优化选访问路径和算法。基于规则：选择下推、投影下推、连接重排、合并选择、拆分条件、消除冗余子查询。基于代价：估计基数、选择率、IO、CPU、内存，比较候选计划。

等价变换：选择串接可交换 $\sigma_{F1}(\sigma_{F2}(R))=\sigma_{F1\land F2}(R)$；选择尽量下推到基表；投影只保留后续需要列；连接满足交换律、结合律，但外连接重排受限制；笛卡尔积+选择转连接；先做高选择性过滤通常更优。

基数估计：选择率越低过滤越强。常用统计：表行数、页数、列基数、最小最大、直方图、索引高度。误差来源：数据倾斜、列相关、过期统计、多谓词独立性假设。连接顺序决定中间结果大小，优化器常用动态规划或启发式搜索；左深树便于流水线，右深树适合某些哈希连接，灌木树搜索空间大。

物理优化：单表访问可选全表扫描、索引扫描、索引范围、覆盖索引；排序可用索引顺序避免 filesort；GROUP BY 可用索引、排序、哈希；LIMIT 与 ORDER BY 合适索引可提前停止。EXPLAIN 看 type、key、rows、filtered、Extra；关注 ALL、大 rows、Using temporary、Using filesort、Using where、Using index。

调优口诀：少取列、早过滤、连键建索引、避免函数包列、避免隐式转换、LIKE 前缀匹配、OR 可考虑 UNION、分页深时用 seek method、统计信息及时更新、用 EXPLAIN 验证。

\blk{MySQL 速查模板}
\begin{lstlisting}
EXPLAIN SELECT ...;
-- seek pagination
SELECT * FROM t WHERE id > ? ORDER BY id LIMIT 20;
-- upsert
INSERT INTO t(id,a) VALUES(1,'x')
ON DUPLICATE KEY UPDATE a=VALUES(a);
-- safe update pattern
UPDATE t SET col=? WHERE pk=? AND version=?;
\end{lstlisting}

\blk{12 分布式数据库}
集中式问题：单点故障、扩展瓶颈、地域延迟。DDB 是逻辑统一、物理分布的数据集合；DDBMS 管理分布式数据；DDBS=DDB+DDBMS+网络+节点+用户。特点：位置透明、复制透明、分片透明、局部自治、分布式查询和事务、故障恢复。

分布式模式：全局外模式、全局概念模式、分片模式、分配模式、局部模式。透明性：用户不关心数据在哪、如何分片、是否复制。数据独立性在分布式下扩展为逻辑、物理、分布独立。

复制：全复制读快可用性高但更新维护贵；部分复制折中；不复制更新简单但可用性低。同步复制一致性强延迟高；异步复制性能好但可能短暂不一致。分片条件：完备性、可重构性、不相交性。水平分片按行，垂直分片按列且需包含键，混合分片组合。分配策略要权衡本地性、负载均衡、通信代价、可用性。

MySQL 主从复制：主库写 binlog，从库 IO 线程拉取 relay log，SQL 线程重放；用途读写分离、备份、容灾；问题是延迟、冲突、单主写瓶颈。MySQL Cluster/NDB 强调内存分布式、高可用、自动分片。

\blk{NoSQL 和 CAP}
关系数据库优势：SQL、事务、完整性、复杂查询、成熟生态；局限：固定模式、水平扩展成本、海量高并发和半结构化数据压力。NoSQL 优势：模式灵活、水平扩展、高吞吐、高可用；代价：弱事务、弱 JOIN、查询能力和一致性取舍。

类型：键值库 Redis/Dynamo 风格，按 key 访问，缓存和会话；文档库 MongoDB，JSON/BSON，内容管理和半结构化；列族库 HBase/Cassandra，宽表稀疏列，海量写入；图数据库 Neo4j，节点边路径，关系分析；向量数据库，向量相似检索，AI 语义搜索。

CAP：Consistency 一致性，Availability 可用性，Partition tolerance 分区容忍。分布式系统必须面对 P；发生网络分区时只能在 C 和 A 之间取舍。CP 保一致牺牲部分可用，AP 保可用接受最终一致。BASE：Basically Available、Soft state、Eventually consistent，与 ACID 相对。

\blk{考场判题和易错清单}
DBMS 不是数据库，DBS 不只是软件。外模式面向用户，内模式面向存储。逻辑独立性由外模式/模式映像保证，物理独立性由模式/内模式映像保证。候选码必须最小，超码不一定最小。主属性属于某候选码，不等于主码属性。自然连接会合并同名列并按同名列等值。投影在关系代数中去重，SQL SELECT 默认不去重。

2NF 只讨论非主属性对候选码的部分依赖，单属性候选码必满足 2NF。3NF 允许右部是主属性，BCNF 不允许非超码决定任何非平凡属性。无损连接不等于保持依赖。E-R 的 m:n 一定转独立关系；1:n 通常在 n 方放外键。

WHERE 不能写聚合条件，聚合条件放 HAVING。NULL 与任何普通比较结果 UNKNOWN。COUNT(col) 忽略 NULL。LEFT JOIN 的右表过滤条件若放 WHERE，可能把外连接变内连接。UNION 去重，UNION ALL 不去重。TRUNCATE 不是带 WHERE 的 DELETE。

复合索引按最左前缀使用；范围条件后列通常不能继续定位；函数包列、隐式转换、前导通配会削弱索引。索引不是越多越好，写多表要控制。视图不等于复制表，可更新性受限制。触发器隐式执行，调试和副作用要谨慎。

事务隔离越高并发越低。脏读、不可重复读、幻读要分清对象：未提交值、同一行值变化、满足条件的行集合变化。行锁依赖索引访问路径，无索引可能扩大锁。死锁不是无限等待，DBMS 可检测并回滚牺牲者。

查询优化不是只靠 SQL 写法，统计信息和数据分布会影响计划。EXPLAIN 的 rows 是估计不是实际。逻辑优化保证等价，物理优化选择算法。外部排序的 $M$ 是缓冲页数，二趟排序条件是 $B\le M(M-1)$。CAP 不是三选二常态，而是分区发生时 C/A 取舍。

\blk{最后 3 分钟检查}
关系题：先列属性和依赖，求 $X^+$ 找候选码，再判 1NF/2NF/3NF/BCNF，分解时分别检查无损和依赖保持。E-R 题：先实体和键，再联系基数，再处理多值属性和 m:n。SQL 题：先 FROM/JOIN，再 WHERE，再 GROUP/HAVING，再 SELECT/ORDER/LIMIT；外连接过滤条件小心位置。索引题：写出查询谓词、排序分组、复合索引顺序，说明最左匹配和代价。事务题：先标出读写冲突和提交点，再判断异常、隔离级别、锁。优化题：写选择下推、投影下推、连接顺序、访问路径和算法代价。分布式题：复制提高读和可用，分片提高容量和并行；CAP 中 P 是现实约束，C/A 是故障时选择。
\end{multicols*}
\end{document}
